\begin{lemma}{Stetigkeit des Linearkombinators}
             {StetigkeitDesLinearkombinators}
Seien $(E, \norm\cdot_E)$ ein normierter $\K$--Vektorraum und $n \in \N$.\\
Wir statten $\K^n$ mit $\norm\cdot_{2,n}$ und $E^n$ mit $\norm\cdot_{2,n,E}$
aus.\\
Dann gilt \[
\forall\, t \in \K^n \forall\, v \in E^n : 
   \norm{\bprod t v _n}_E \le \norm t_{2, n} \cdot \norm v _{2, n, E}
\]
Insbesondere ist $\bprod\cdot{\cdot\cdot}_n$ (mit genannter Normausstattung)
stetig.
\end{lemma}
\begin{proof}
Seien $t \in \K^n$ und $v \in E^n$.
Dann gilt
\begin{align*}
\norma{\bprod t v_n}_E 
    &= \norma{\sum_{i=1}^n t_iv_i}_E 
    \le\sum_{i=1}^n \abs{t_i}\cdot \norm{v_i}_E\\
    &= \sprod{\left(\pfeil {\abs\cdot}n(t)\right)}
            {\left(\pfeil {\norm\cdot_E}n(v)\right)}_n\\
    &\le\norma{\pfeil {\abs\cdot}n(t)}_{2,n}
        \cdot
        \norma{\pfeil{\norm\cdot_E}n(v)}_{2, n} & \quad\text{Cauchy--Schwarz}\\
    &= \norm t_{2, n}\cdot \norm v_{2, n, E}
\end{align*}
Da $\bprod\cdot{\cdot\cdot}_n$ nach \ref{EigenschaftenDesLinearkombinators}
bilinear ist, folgt mit der bekannten Charakterisierung der Stetigkeit
bilinearer Abbildungen die Stetigkeit für den Linearkombinator.
\end{proof}